\subsection{Modelos aprobados}

De los cinco enfoques significativos que reportamos en la etapa de Modelado, solo el Modelo 5 con salida redondeada superó conjuntamente las metas de error y tolerancia. Es importante distinguir entre aquellos modelos que no pudieron alcanzar los criterios de éxito y aquellos cuyo entrenamiento no terminó. En particular, el Modelo 3 (EVA02 con ajuste fino) no cuenta con resultados suficientes como para realizar una comparación válida, de modo que lo consideramos como no evaluable en esta fase y no puede aprobarse ni descartarse por su desempeño predictivo.

\begin{table}[H]
  \centering
  \small
  \caption{Resultados de los modelos significativos sobre las 76 vacas.}
  \label{tab:modelos}
  \begin{tabularx}{\textwidth}{>{\raggedright\arraybackslash}X>{\raggedright\arraybackslash}p{0.08\textwidth}>{\raggedright\arraybackslash}p{0.13\textwidth}>{\raggedright\arraybackslash}p{0.08\textwidth}>{\raggedright\arraybackslash}p{0.15\textwidth}}
    \toprule
    Modelo & MAE & $\pm0.25$ continua & $\pm0.50$ & $\pm0.25$ redondeada \\
    \midrule
    1. CNN propia & 0.341 & 42.1\,\% & 77.6\,\% & 52.6\,\% \\
    \addlinespace
    2. ConvNeXt-Tiny & 0.255 & 63.2\,\% & 86.8\,\% & 73.7\,\% \\
    \addlinespace
    3. EVA02 & --- & --- & --- & --- \\
    \addlinespace
    4. ResNet-50 & 0.260 & 61.8\,\% & 85.5\,\% & 75.0\,\% \\
    \addlinespace
    5. Clasificación con valor esperado & 0.236 & 64.5\,\% & 90.8\,\% & 77.6\,\% \\
    \addlinespace
    5. Salida final redondeada & 0.224 & --- & 94.7\,\% & 77.6\,\% \\
    \bottomrule
  \end{tabularx}
\end{table}

Las cifras de $\pm0.50$ para los modelos 1 a 5 antes del redondeo son de su salida continua, mientras que el 94.7\,\% del modelo final es de la salida redondeada. Separamos estas modalidades para evitar tener comparaciones que sean engañosas.

\paragraph{Condiciones de aprobación técnica.}

\begin{enumerate}
  \renewcommand{\labelenumi}{(\roman{enumi})}
  \item Conservar la salida en pasos de 0.25 y documentar la transformación.
  \item Reportar conjuntamente el rendimiento continuo y el rendimiento redondeado.
  \item Usarlo solo como apoyo que estará sujeto a revisión por parte del experto.
  \item Advertir sobre el sesgo hacia 2.75--3.25 y tratar las estimaciones en extremos con cuidado.
  \item Validar a futuro la captura, consistencia y errores por rango antes de realizar un despliegue autónomo.
\end{enumerate}
